% Source-aligned English for printed p.140, with logged mathematical correction.
\SourcePage{140}

\LayerTarget{tohoku:axiom:delta-functor:2:p2}
\noindent
\[
\tag{2.1.1.}
\ldots\,.\;T^i(A')\to T^i(A)
\to T^i(A'')\to T^{i+1}(A')
\to\ldots\,.
\]

\noindent is a complex; that is, the composite of any two consecutive
morphisms in this sequence is zero.

\LayerTarget{tohoku:def:delta-star-and-variance}
The definition of a covariant $\partial^*$-functor is analogous, the only
difference being that the operator $\partial$ lowers degree by one instead of
raising it. The definitions of contravariant $\partial$-functors and
$\partial^*$-functors are likewise analogous: the $T^i$ are then
contravariant additive functors, and the boundary operators go from
$T^i(A')$ to $T^{i+1}(A'')$ or $T^{i-1}(A'')$. Reversing the signs of the
degrees $i$ in $T^i$, or replacing $\mathbf{C}'$ by its opposite category, transforms
$\partial$-functors into $\partial^*$-functors. Thus one may always reduce to
the study of covariant $\partial$-functors. Observe that when $a=-\infty$ and
$b=+\infty$, a $\partial$-functor is a \textquotedblleft connected sequence
of functors\textquotedblright\ in [6, Chapter III].

\LayerTarget{tohoku:def:delta-functor-morphism}
Given two $\partial$-functors $T$ and $T'$ defined in the same degrees, a
\textit{morphism}, or natural transformation, from $T$ to $T'$ is a system
$f=(f^i)$ of natural transformations $f^i:T^i\to T'^i$ satisfying
the natural compatibility condition with $\partial$: for every exact
sequence
$0\to A'\to A\to A''\to 0$,
the diagram

\LayerTarget{tohoku:diagram:delta-morphism}
\[
\begin{array}{ccc}
T^i(A'') & \mathop{\to}\limits^{\partial} & T^{i+1}(A')\\
\big\downarrow && \big\downarrow\\
T'^i(A'') & \mathop{\to}\limits^{\partial} & T'^{i+1}(A')
\end{array}
\]

\noindent is commutative. Morphisms of $\partial$-functors add and compose in
the evident way.

\LayerTarget{tohoku:def:exact-delta-functor}
Suppose that $\mathbf{C}'$ is also an abelian category. A $\partial$-functor is called
\textit{exact} if, for every exact sequence
$0\to A'\to A\to A''\to 0$
in $\mathbf{C}$, the corresponding sequence (2.1.1) is exact. A
\textit{cohomological functor}, respectively a \textit{homological functor},
is an \textit{exact} $\partial$-functor, respectively $\partial^*$-functor, defined in
all degrees.

\LayerTarget{tohoku:section:2.2}
\indent\textbf{2.2. Universal $\partial$-functors.}
\LayerTarget{tohoku:def:universal-delta-functor}
Let $T=(T^i)\quad 0\leqq i\leqq a$ be a covariant $\partial$-functor from $\mathbf{C}$ to
$\mathbf{C}'$, where $a>0$. The functor $T$ is called \textit{universal} if, for every
$\partial$-functor $T'=(T'^i)$ defined in the same degrees and every natural
transformation $f^0:T^0\to T'^0$, there exists a
unique morphism $f:T\to T'$ whose component in degree $0$ is
$f^0$. The definition is identical for a contravariant
$\partial$-functor; for $\partial^*$-functors one must instead consider
morphisms from $T'$ to $T$, rather than from $T$ to $T'$.

\LayerTarget{tohoku:prop:satellites-uniqueness}
By the definition itself, given a covariant functor $F:\mathbf{C}\to \mathbf{C}'$
and an integer $a>0$, there can exist, up to canonical isomorphism, at most
one universal $\partial$-functor defined in degrees $0\leqq i\leqq a$ and equal
to $F$ in degree zero.
\LayerTarget{tohoku:def:right-left-satellites}
Its components are denoted $S^iF$ and called the right \textit{satellite
functors} of $F$. For $i>0$ we also set $S_iF=S^{-i}F$, where the $S_iF$ are
the \textit{left} satellites of $F$, defined like the $S^iF$ but by using
universal $\partial^*$-functors in degrees $0\leqq i\leqq a$ such that $T^0=F$.
It follows immediately that, for fixed $i$, if $S^iF$ exists, it is
independent of the choice of $a$.

\LayerTarget{tohoku:remark:satellites-existence:p1}
In every case known to me, the satellite functors of an arbitrary additive
functor $F$ exist. Moreover, given $\mathbf{C}$ and $\mathbf{C}'$, in order to prove

\TohokuContinuation{141}
